\scanpage{17}\mspage{6}
La catégorie des objets admissibles de $\widehat C$ est stable par \deleted{\unsure{les} l'opération} sommes quelconques, et par sommes amalgamées $X\amalg_Z Y$, \unsure{don|dont} $Z\unsure{\mathord{-}|\to}X$ et $Z\mathord{-}Y$ des mono, et le foncteur $X\to\widetilde X$ commute : ces opérations, dans un sens évident. De plus, les objets de type fini donnent des objets de t.f., pourvu, \unsure{(}dans le cas des sommes\unsure{)|,}, qu'on se borne \unsure{:|à} des sommes \uline{finies}.

\uline{NB} Quand les $\unsure{J_n|I_n|J}$ sont finis, alors les objets \inserted{adm.} de t.f. de $\widehat C$ \deleted{\unsure{aussi|sont} \unsure{(par ex)|t.f. les $X$ t.|de les $X$,}} \inserted{correspondent aux} ens ordonnés $M$-épinglés $K$ qui \deleted{sont} \uline{finis}.
% archive sheet number: 16
